\documentclass[../main.tex]{subfiles}
\begin{document}

\chapter{Predication}
\lessoninfo{
  video={Lesson 5, videos 1--12},
  textbook={H\&P \cite{hennessy2017} §A.6, §C.2, §H.4, §H.6},
  keywords={predication, hard-to-predict branch, if-conversion, conditional move, MOVZ, MOVN, CMOVcc, condition code, break-even misprediction rate, full predication, Itanium, predicate register},
}

The predictors of Lesson 4 guess most branches correctly, but some branches
remain hard to predict for any predictor, typically because their direction
depends on the data being processed, and in a deep pipeline every
misprediction flushes tens of instructions.  This lesson studies the
alternative for such branches: \emph{predication}, which executes the work of
both directions of a branch and keeps only the results of the correct one, so
that nothing is left to mispredict.  We examine which branches benefit, how a
compiler removes a branch by if-conversion, the conditional move instructions
this relies on, when the transformation pays off, and instruction sets in
which every instruction can be made conditional.

\section{Prediction and predication}
\label{sec:l05-predication}

A branch predictor guesses in the fetch stage where a branch goes, and the
processor continues to fetch along the guessed path.\source{Lesson 5, videos
1--2; H\&P §C.2.}  A correct guess costs nothing.  A wrong guess flushes every
instruction fetched after the branch, which in a modern processor, several
instructions wide and many stages deep, amounts to tens of instructions
(\cref{ex:l04-better}).  The cost of a branch is thus zero most of the time
and large once in a while.

\begin{definition}[Predication]
  \label{def:l05-predication}
  Executing the work of \emph{both} directions of a branch and throwing away
  the results of the direction that turns out to be wrong is called
  \term{predication}[プレディケーション].
\end{definition}

Predication has no misprediction penalty, since nothing is predicted, but it
always wastes the work of the direction that is not needed: if both
directions contain equally many instructions, half of the work is thrown
away.\caution{Nothing is guessed and nothing is undone: both directions run
on purpose, and the code itself discards the unneeded results.}  Per
execution of the branch, predication wastes the instructions of the unneeded
direction, while prediction wastes the misprediction penalty
multiplied by the probability of a misprediction.  Which is smaller depends
on the kind of branch (\cref{tab:l05-which}); predication can be the better
choice only for a small if-then-else, whose directions contain far fewer
instructions than a misprediction flushes.

\begin{table}[htbp]
  \centering\small
  \caption{Prediction and predication for different kinds of branches.}
  \label{tab:l05-which}
  \begin{tabular}{@{}>{\raggedright\arraybackslash}p{0.17\linewidth}>{\raggedright\arraybackslash}p{0.27\linewidth}>{\raggedright\arraybackslash}p{0.27\linewidth}>{\raggedright\arraybackslash}p{0.19\linewidth}@{}}
    \toprule
    Branch & Prediction & Predication & Better \\
    \midrule
    Loop               & accurate
                       & also runs the code after the loop in every iteration
                       & prediction \\
    \addlinespace
    Call, return       & always taken; BTB and RAS give the target
                       & work on the other path is always wasted
                       & prediction \\
    \addlinespace
    Large if-then-else & occasionally flushes tens of instructions
                       & wastes a large direction every time
                       & prediction \\
    \addlinespace
    Small if-then-else & flushes tens of instructions per misprediction
                       & wastes a few instructions every time
                       & predication, if hard to predict \\
    \bottomrule
  \end{tabular}
\end{table}

For a small if-then-else the choice depends on how well the branch is
predicted.  If a misprediction flushes 40 instructions and the unneeded
direction contains 4, predication wastes less than prediction as soon as more
than one execution in ten is mispredicted.\margin{Real penalties are
larger: about 17 cycles on a Skylake-class core (Lesson 4), some 70
instruction slots at four instructions per cycle.}
\Cref{sec:l05-performance} makes this comparison precise for the code that a
compiler actually produces.

\begin{keypoint}
  Prediction costs nothing when it is right and tens of instructions when it
  is wrong; predication always wastes the unneeded direction.  Branches are
  therefore predicted by default, and predication is worth considering only
  for small if-then-else branches that are hard to predict.
\end{keypoint}

\section{If-conversion}
\label{sec:l05-if-conversion}

Predication needs the work of both directions in the form of instructions
that are executed one after the other, with no branch between
them.\source{Lesson 5, video 3; H\&P §H.4; \cite{allen1983}.}  Producing
such code is the task of the compiler.  Consider the following if-then-else, which sums the
values below and above a threshold separately and counts the values below it.
If the values of \texttt{v} fall on either side of \texttt{t} in no regular
order, the branch is hard to predict.
\begin{lstlisting}[language=C, numbers=none]
if (v < t) {
  lo = lo + v;
  n  = n + 1;
} else {
  hi = hi + v;
}
\end{lstlisting}

\begin{definition}[If-conversion]
  \label{def:l05-if-conversion}
  The compiler transformation that removes such a branch by turning the
  control dependence (Lesson 3) of the statements of both directions into a
  data dependence on the branch condition is called
  \term{if-conversion}[if 変換].
\end{definition}

With ordinary instructions, the converted code (\cref{fig:l05-if-conversion})
executes the statements of both directions unconditionally, writing their
results into temporaries, and then uses the condition~\texttt{c} as data to
select the new value of every variable assigned in either
direction.\caution{The unneeded direction also runs: it must not store
to memory or fault, e.g.\ by loading via an invalid pointer.}  When both
directions assign a variable~\texttt{x}, the selection chooses between their
two temporaries, \texttt{x = c ? x1 : x2}; when only one direction assigns
it, as for each variable in this example, the other alternative is the old
value.  The converted code contains no branch.

\begin{figure}[htbp]
  \centering
  % If-conversion: a small if-then-else (control flow) becomes a straight
% sequence that computes both parts and selects the results (data flow).
\begin{tikzpicture}[x=1cm,y=1cm, font=\footnotesize\sffamily, >={Stealth[length=1.8mm]},
  blk/.style={draw, fill=ln-box, align=left, inner sep=3pt, font=\footnotesize\ttfamily},
  sel/.style={blk, fill=ln-accent!15},
  lab/.style={font=\scriptsize\sffamily, text=ln-muted}]
  % ---- (a) control flow
  \node[blk] (c) at (2.1,0) {c = (v < t)};
  \node[draw, diamond, aspect=1.6, inner sep=1pt, fill=white] (br) at (2.1,-1.05) {c?};
  \node[blk] (th) at (1.0,-2.45) {lo = lo + v\\n\ \ = n + 1};
  \node[blk] (el) at (3.3,-2.45) {hi = hi + v};
  \node[circle, draw, fill=white, inner sep=1.3pt] (j) at (2.1,-3.65) {};
  \draw[->] (c) -- (br);
  \draw[->] (br.west) -| node[pos=0.25, above, font=\scriptsize\sffamily] {T} (th.north);
  \draw[->] (br.east) -| node[pos=0.25, above, font=\scriptsize\sffamily] {N} (el.north);
  \draw[->] (th.south) |- (j.west);
  \draw[->] (el.south) |- (j.east);
  \draw[->] (j.south) -- ++(0,-0.35);
  \node[font=\footnotesize\sffamily] at (2.1,-4.55) {\iflangja{(a) 分岐}{(a) branch}};
  % ---- transformation arrow
  \draw[-{Stealth[length=2.6mm]}, very thick, ln-accent] (4.45,-1.8) -- (5.55,-1.8)
    node[midway, above, font=\scriptsize\sffamily, text=ln-accent, align=center]
    {\iflangja{if 変換}{if-\\conversion}};
  % ---- (b) data flow
  \node[blk, anchor=north west, minimum width=3.3cm] (c2) at (5.9,0.3) {c = (v < t)};
  \node[blk, anchor=north west, minimum width=3.3cm] (t2) at (5.9,-0.45) {lo1 = lo + v\\n1\ \ = n + 1};
  \node[blk, anchor=north west, minimum width=3.3cm] (e2) at (5.9,-1.55) {hi1 = hi + v};
  \node[sel, anchor=north west, minimum width=3.3cm] (s2) at (5.9,-2.3) {lo = c ? lo1 : lo\\n\ \ = c ? n1\ \ : n\\hi = c ? hi\ \ : hi1};
  \draw[->] (c2) -- (t2); \draw[->] (t2) -- (e2); \draw[->] (e2) -- (s2);
  \draw[->, dashed, ln-accent] (c2.east) -- ++(0.25,0) |- (s2.east);
  \node[lab, anchor=west] at (9.5,-0.75) {\iflangja{then 部}{then part}};
  \node[lab, anchor=west] at (9.5,-1.7) {\iflangja{else 部}{else part}};
  \node[lab, anchor=west] at (9.5,-2.9) {\iflangja{選択}{select}};
  \draw[->] (s2.south) -- ++(0,-0.35);
  \node[font=\footnotesize\sffamily] at (7.55,-4.55) {\iflangja{(b) if 変換後}{(b) if-converted}};
\end{tikzpicture}

  \caption{If-conversion of a small if-then-else.  (a) The original code
    branches on~\texttt{c}.  (b) The converted code computes both directions
    into temporaries and uses~\texttt{c} as data to select the results.}
  \label{fig:l05-if-conversion}
\end{figure}

The transformation gains something only if the selections can be carried out
without a branch.\margin{Removing the branch also leaves one straight-line
stretch of code, which the compiler can then schedule as a whole (Lesson
10).}  Translated into ordinary branch instructions, each
selection is itself an if-then-else on~\texttt{c}, as hard to predict as the
branch that was removed, so the converted code would still mispredict and in
addition execute the work of both directions.  Ordinary arithmetic can select
without a branch by masking both alternatives with AND and OR, but that costs
several instructions per selection; if-conversion is cheap only if the
instruction set can make the effect of an instruction conditional.

\section{Conditional move instructions}
\label{sec:l05-cmov}

The simplest such instruction makes only a copy conditional.\source{Lesson 5,
videos 4--5; H\&P §H.4, §A.6.}

\begin{definition}[Conditional move]
  \label{def:l05-cmov}
  A \term{conditional move}[条件付きムーブ] is an instruction that copies a
  value into its destination if a condition holds and leaves the destination
  unchanged otherwise.
\end{definition}

The instruction is executed in either case; only whether it writes depends on
the condition.  It does not change the flow of control, so there is nothing to
predict.

The MIPS instruction set has two conditional moves whose condition is the
contents of a register.  The instruction \term{MOVZ}[MOVZ] (move if zero)
copies register Rs into register Rd if register Rt holds zero, and the
instruction \term{MOVN}[MOVN] (move if not zero) copies it if Rt does not hold
zero:
\begin{lstlisting}[language={[x86masm]Assembler}, numbers=none]
MOVZ Rd, Rs, Rt     ; if (Rt == 0) Rd <- Rs
MOVN Rd, Rs, Rt     ; if (Rt != 0) Rd <- Rs
\end{lstlisting}
A selection \texttt{x = c ? x1 : x2} becomes two moves into the register of
\texttt{x}, of which exactly one writes.  With \texttt{c} in R3, \texttt{x1}
and \texttt{x2} in R1 and R2, and \texttt{x} in R4:
\begin{lstlisting}[language={[x86masm]Assembler}, numbers=none]
MOVN R4, R1, R3     ; if c != 0: x <- x1
MOVZ R4, R2, R3     ; if c == 0: x <- x2
\end{lstlisting}
The condition register is typically set by a comparison such as
\texttt{SLT}, which writes 1 if the condition holds and 0
otherwise.\tip{After \texttt{SLT} the condition register is nonzero exactly
when the condition holds: the \texttt{MOVN} carries the then part and the
\texttt{MOVZ} the else part.}

\begin{example}
  \label{ex:l05-movzn}
  With \texttt{v} in R1, \texttt{t} in R2, and \texttt{lo}, \texttt{n},
  \texttt{hi} in R4, R5, R6, the code of \cref{sec:l05-if-conversion}
  compiles with a branch into
  \begin{lstlisting}[language={[x86masm]Assembler}, numbers=none]
      SLT  R3, R1, R2     ; R3 <- (v < t)
      BEQ  R3, R0, Skip   ; v >= t: go to else part
      ADD  R4, R4, R1     ; lo <- lo + v
      ADDI R5, R5, 1      ; n  <- n + 1
      J    Cont           ; skip else part
Skip: ADD  R6, R6, R1     ; hi <- hi + v
Cont: ...
\end{lstlisting}
  and, if-converted, into
  \begin{lstlisting}[language={[x86masm]Assembler}, numbers=none]
      SLT  R3, R1, R2     ; R3 <- (v < t)
      ADD  R7, R4, R1     ; then part into R7, R8
      ADDI R8, R5, 1
      ADD  R9, R6, R1     ; else part into R9
      MOVN R4, R7, R3     ; lo <- R7 if v < t
      MOVN R5, R8, R3     ; n  <- R8 if v < t
      MOVZ R6, R9, R3     ; hi <- R9 if v >= t
\end{lstlisting}
  The temporaries R7--R9 are needed because R4--R6 must keep their old
  values until the moves decide which of them change.  Each variable here is
  assigned in only one direction, so a single move per variable suffices.
\end{example}

Other instruction sets base conditional moves on flags rather than on the
contents of a register.

\begin{definition}[Condition code]
  \label{def:l05-condition-code}
  A \term{condition code}[条件コード], or flag, is a one-bit status value, for
  example zero, negative, carry or overflow, that arithmetic and comparison
  instructions set as a side effect and that later instructions test.
\end{definition}

The x86 instruction set has a family of conditional moves,
\term{CMOVcc}[CMOVcc], in which \texttt{cc} names a condition on the condition
codes: \texttt{CMOVZ} moves if the zero flag is set, \texttt{CMOVNZ} if it is
clear, \texttt{CMOVG} if the last comparison found \enquote{greater},
\texttt{CMOVL} if it found \enquote{less}, and so on.  A common pattern
assigns one alternative unconditionally and lets a single conditional move
overwrite it with the other, as in this computation of $m = \max(a,b)$:
\begin{lstlisting}[language={[x86masm]Assembler}, numbers=none]
MOV   EAX, EBX      ; m <- a
CMP   EBX, ECX      ; set flags from a - b
CMOVL EAX, ECX      ; if a < b: m <- b
\end{lstlisting}

\section{Performance of if-conversion}
\label{sec:l05-performance}

If-conversion is not free: the converted code executes the work of both
directions and the conditional moves on every execution.\source{Lesson 5,
videos 6--7; H\&P §C.2.}  To compare the two versions of a code fragment we
count pipeline slots, as in Lesson 3: every executed instruction occupies one
slot, and every misprediction wastes as many slots as it flushes instructions.
Let the branch version execute $n_b$ instructions per execution on average
and mispredict $m$ times per execution with a penalty of $P$ instructions, and
let the converted version execute $n_c$ instructions.\caution{$m$ is the
rate of this one branch, per execution of the fragment, not the
per-instruction rate of Lesson 4; $P$ counts instructions, not cycles.}
If-conversion pays off when
\begin{equation}
  n_c \;<\; n_b + mP
  \qquad\Longleftrightarrow\qquad
  m \;>\; \frac{n_c - n_b}{P} .
  \label{eq:l05-break-even}
\end{equation}
The right-hand side is the \emph{break-even misprediction rate}: the extra
instructions of the converted code, measured in units of the misprediction
penalty.\tip{Each extra instruction of the converted code raises the
break-even rate by $1/P$: with $P = 20$, five percentage points of predictor
accuracy apiece.}  Compared with the rough comparison of
\cref{sec:l05-predication}, the extra instructions include the conditional
moves, less the branch and jump instructions that the conversion removes.

\begin{example}
  \label{ex:l05-performance}
  The branch version of \cref{ex:l05-movzn} executes five instructions when
  the then part applies (\texttt{SLT}, \texttt{BEQ}, \texttt{ADD},
  \texttt{ADDI}, \texttt{J}) and three otherwise (\texttt{SLT}, \texttt{BEQ},
  \texttt{ADD}).  If both parts apply equally often, $n_b = 4$; the converted
  version always executes $n_c = 7$.  Suppose the predictor is
  \SI{75}{\percent} accurate on this branch and $P = 20$.  The branch version
  costs $4 + 0.25 \times 20 = 9$ slots per execution and the converted
  version 7, a speedup of $9/7 \approx 1.29$ for the fragment.  By
  \cref{eq:l05-break-even} the break-even rate is $m = 3/20 = 0.15$: the
  branch version is faster only if the predictor is more than
  \SI{85}{\percent} accurate.  With $P = 40$ the break-even rate halves to
  $0.075$, an accuracy of \SI{92.5}{\percent}.
\end{example}

\begin{figure}[htbp]
  \centering
  % Slots per execution of the example fragment as a function of the
% misprediction rate: branch version (4 instructions on average, penalties
% 20 and 40), MOVZ/MOVN version (7 instructions), fully predicated version
% (4 instructions).
\begin{tikzpicture}[x=26cm, y=0.42cm, font=\footnotesize\sffamily, >={Stealth[length=1.8mm]}]
  % grid and axes
  \foreach \yv in {2,4,6,8,10} { \draw[ln-rule!50, very thin] (0,\yv) -- (0.3,\yv); }
  \draw[->] (0,0) -- (0.315,0);
  \node[font=\scriptsize\sffamily] at (0.15,-1.75)
    {\iflangja{実行当たり予測ミス数 $m$}{mispredictions per execution, $m$}};
  \draw[->] (0,0) -- (0,10.8) node[above, font=\scriptsize\sffamily]
    {\iflangja{実行当たりのスロット数}{slots per execution}};
  \foreach \xv/\xl in {0/0,0.05/0.05,0.1/0.10,0.15/0.15,0.2/0.20,0.25/0.25,0.3/0.30} {
    \draw (\xv,0) -- (\xv,-0.2) node[below, font=\scriptsize\sffamily] {\xl};
  }
  \foreach \yv in {0,2,4,6,8,10} {
    \draw (0,\yv) -- (-0.004,\yv) node[left, font=\scriptsize\sffamily] {\yv};
  }
  % break-even guides (drawn first, under the lines)
  \foreach \xb/\yb in {0.075/7,0.15/7} {
    \draw[ln-muted, densely dotted] (\xb,\yb) -- (\xb,0);
  }
  % lines
  \draw[very thick, ln-caution] (0,4) -- (0.3,4)
    node[pos=0.8, below, font=\scriptsize\sffamily, text=ln-caution] {\iflangja{完全プレディケーション}{full predication}};
  \draw[very thick, ln-tip] (0,7) -- (0.3,7)
    node[pos=0.82, below, font=\scriptsize\sffamily, text=ln-tip] {MOVZ/MOVN};
  \draw[very thick, ln-accent] (0,4) -- (0.3,10)
    node[pos=0.8, above left, font=\scriptsize\sffamily, text=ln-accent] {\iflangja{分岐，$P=20$}{branch, $P=20$}};
  \draw[very thick, ln-accent, dashed] (0,4) -- (0.15,10)
    node[pos=0.85, left, font=\scriptsize\sffamily, text=ln-accent, xshift=-1mm] {\iflangja{分岐，$P=40$}{branch, $P=40$}};
  % break-even points
  \foreach \xb/\yb in {0.075/7,0.15/7,0/4} {
    \fill (\xb,\yb) circle (1.4pt);
  }
\end{tikzpicture}

  \caption{Slots per execution of the fragment of \cref{ex:l05-performance}
    as a function of the misprediction rate.  Dots mark the break-even
    rates; the fully predicated version is discussed in
    \cref{sec:l05-full}.}
  \label{fig:l05-break-even}
\end{figure}

\Cref{fig:l05-break-even} shows the comparison as a function of~$m$.  The
cost of the branch version grows linearly with the misprediction rate, with a
slope equal to the penalty, while the converted version costs the same for
every~$m$.  A deeper or wider pipeline raises the penalty and so favours
if-conversion; larger then and else parts raise $n_c - n_b$ and favour the
branch.

\begin{keypoint}
  If-conversion replaces an occasional large penalty by a constant overhead.
  It pays off when the mispredictions per execution times the penalty exceed
  the extra instructions of the converted code (\cref{eq:l05-break-even}).
\end{keypoint}

\needspace{8\baselineskip}
\section{Benefits and costs of conditional moves}
\label{sec:l05-cmov-summary}

Removing branches with conditional moves has one benefit and several
costs.\source{Lesson 5, video 8; H\&P §H.4.}
\begin{itemize}
  \item \textbf{Benefit:} hard-to-predict branches, and with them their
    mispredictions, disappear.
  \item \textbf{Compiler support} is needed.  The processor executes the code
    it is given; a branch is removed only if the compiler recognizes a
    suitable if-then-else and converts it.
  \item \textbf{More registers} are needed: every result is first computed
    into a temporary.
  \item \textbf{More instructions} are executed: the work of both directions,
    even on executions where the branch would have been predicted correctly,
    and the conditional moves that select the results.
\end{itemize}

Of these, the compiler support, the removal of the branch and the execution
of both directions are inherent to predication; the extra registers and the
moves are not.  They arise because only the move is conditional: every
result has to be computed into a temporary so that a conditional move can
decide whether to keep it.  If the instruction
that computes a result could itself be made conditional, neither the
temporaries nor the moves would be needed.

\section{Full predication}
\label{sec:l05-full}

A conditional move is a separate opcode, and the only operation it makes
conditional is a copy.\source{Lesson 5, videos 9--12; H\&P §H.4, §H.6;
\cite{mahlke1995}.}
The alternative is support in the instruction set for making every
instruction conditional.

\begin{definition}[Full predication]
  \label{def:l05-full-predication}
  An instruction set provides \term{full predication}[完全プレディケーション]
  if condition bits are added to every instruction: the instruction takes
  effect only if its condition holds, and has no effect otherwise.
\end{definition}

The IA-64 instruction set of Intel's Itanium processors is an
example\index{Itanium} (\cref{fig:l05-qp}).  In most IA-64 instructions, the
six least significant bits of the 41-bit instruction hold a \emph{qualifying
predicate}: the number of one of 64 one-bit registers, each called a
\term{predicate register}[述語レジスタ].  The instruction is fetched and
occupies its pipeline slot in either case, but it takes effect only if the
selected predicate register holds~1; otherwise it is annulled: it writes no
result and raises no exception.  Predicate register~0 always holds~1, so an
instruction whose qualifying predicate is~0 is unconditional.
The predicate registers are set by compare instructions, which write a pair
of predicates with complementary values, one for the condition and one for
its negation.\margin{A 128-bit bundle holds three 41-bit instructions and a
5-bit template; IA-64 also has 128 integer and 128 floating-point registers
(Lesson 11).}

\begin{figure}[htbp]
  \centering
  % Full predication in IA-64: the 6-bit qualifying predicate (qp) of every
% 41-bit instruction selects one of 64 one-bit predicate registers, and the
% instruction runs in either case, but its result is written only if that
% register holds 1; otherwise it is annulled (no write, no exception).
\begin{tikzpicture}[x=1cm,y=1cm, font=\footnotesize\sffamily, >={Stealth[length=1.8mm]},
  box/.style={draw, fill=ln-box, align=center, minimum height=0.8cm, inner sep=3pt},
  pr/.style={draw, fill=white, minimum width=0.42cm, minimum height=0.42cm, inner sep=0pt, font=\scriptsize\ttfamily},
  lab/.style={font=\scriptsize\sffamily, text=ln-muted}]
  % ---- instruction word, bits 40..0
  \draw[fill=ln-box] (0,0) rectangle (7.0,0.6);
  \draw[fill=ln-accent!15] (7.0,0) rectangle (8.4,0.6);
  \node at (3.5,0.3) {\iflangja{オペコードとオペランド}{opcode and operands}};
  \node at (7.7,0.3) {qp};
  \foreach \x/\b in {0.15/40,6.85/6,7.15/5,8.25/0} {
    \node[lab, above] at (\x,0.6) {\b};
  }
  % ---- operation and write
  \node[box, minimum width=3.3cm] (op) at (1.75,-1.3) {\iflangja{演算を実行}{perform the operation}};
  \node[box, minimum width=3.3cm] (wr) at (1.75,-3.0)
    {\iflangja{結果を書き込む}{write the result}\\[-1pt]
     {\scriptsize\iflangja{（レジスタまたはメモリ）}{(register or memory)}}};
  \draw[->] (1.75,0) -- (op.north);
  \draw[->] (op.south) -- (wr.north);
  \node[draw, dashed, ln-caution, inner sep=5pt, fit=(wr)] (fr) {};
  % ---- predicate registers PR0..PR63 (one bit each), PR3 selected
  \foreach \v [count=\k from 0] in {1,0,1,1,0} {
    \node[pr] (p\k) at (7.7+0.42*\k-1.26,-1.3) {\v};
  }
  \node[pr, draw=none, fill=none] at (7.7+0.42*5-1.26,-1.3) {$\cdots$};
  \node[pr] (p63) at (7.7+0.42*6-1.26,-1.3) {0};
  \node[pr, fill=ln-accent!15] at (p3) {1};
  \node[lab, below] at (p0.south) {0};
  \node[lab, below] at (p63.south) {63};
  \node[lab, left, align=right] at ([xshift=-1mm]p0.west)
    {\iflangja{述語レジスタ}{predicate registers}\\\iflangja{（各 1 ビット）}{(1 bit each)}};
  \node[lab, anchor=north west, align=left] at (8.0,-2.0)
    {\iflangja{比較命令が設定\\0 番は常に 1}{set by compare instructions;\\register 0 always holds 1}};
  % qp selects a register
  \draw[->, ln-accent, thick] (7.7,0) -- (p3.north)
    node[midway, right, font=\scriptsize\sffamily, text=ln-accent] {qp = 3};
  % enable
  \draw[->, ln-caution, thick] (p3.south) |- (fr.east |- wr.east)
    node[pos=0.75, above, font=\scriptsize\sffamily, text=ln-caution]
    {\iflangja{1 のときのみ結果を書き込む}{result written only if the bit is 1}}
    node[pos=0.75, below, font=\scriptsize\sffamily, text=ln-caution, align=center]
    {\iflangja{0 なら無効化\\（書き込みも例外もなし）}{if 0: annulled\\(no write, no exception)}};
\end{tikzpicture}

  \caption{Full predication in IA-64.  The qualifying predicate (qp) of an
    instruction selects one of 64 one-bit predicate registers.  The
    instruction passes through the pipeline in either case, but its result is
    written only if that register holds~1; otherwise the instruction is
    annulled, with no write and no exception.}
  \label{fig:l05-qp}
\end{figure}

\needspace{8\baselineskip}
\begin{example}
  \label{ex:l05-full}
  With full predication, if-conversion guards each statement of both
  directions directly with the condition or its negation, and the code of
  \cref{sec:l05-if-conversion} becomes
  \begin{lstlisting}[language={[x86masm]Assembler}, numbers=none]
      CMPLT P1, P2, R1, R2 ; P1 <- (v < t), P2 <- !P1
(P1)  ADD   R4, R4, R1     ; lo <- lo + v  if v < t
(P1)  ADDI  R5, R5, 1      ; n  <- n + 1   if v < t
(P2)  ADD   R6, R6, R1     ; hi <- hi + v  if v >= t
\end{lstlisting}
  where the qualifying predicate is written in parentheses before the
  instruction and \texttt{CMPLT} denotes a compare instruction that sets P1
  to~1 and P2 to~0 if R1 is less than R2, and the reverse otherwise.  The additions
  themselves are conditional, so they write directly into the variables, and
  the compare replaces both the \texttt{SLT} and the branch.
  \Cref{tab:l05-versions} compares the three versions with the numbers of
  \cref{ex:l05-performance}.
\end{example}

\begin{table}[htbp]
  \centering\small
  \caption{Three versions of the example fragment with a \SI{75}{\percent}
    accurate predictor and $P = 20$.}
  \label{tab:l05-versions}
  \begin{tabular}{@{}lccccc@{}}
    \toprule
    Version & Instructions & Temporaries & Moves & Mispred. & Slots \\
    \midrule
    Branch           & 5 or 3 & 0 & 0 & 0.25 & 9 \\
    MOVZ/MOVN        & 7      & 3 & 3 & 0    & 7 \\
    Full predication & 4      & 0 & 0 & 0    & 4 \\
    \bottomrule
  \end{tabular}
\end{table}

The fully predicated version executes the compare and the instructions of
both directions, and nothing else.\sidenote{IA-64 was developed jointly by HP
and Intel from 1994, and full predication was part of its plan to expose
parallelism at compile time instead of in hardware.  The first Itanium shipped
in 2001; it never displaced x86-64, and Intel shipped the last one in July
2021 (Lesson 11).}  Compared with the branch version it
executes the unneeded direction but no branch and no jump; here the two
balance, and its $n_c = 4$ equals the average $n_b$ of the branch version.
Its break-even rate by \cref{eq:l05-break-even} is $m = (4-4)/20 = 0$: the
branch version is not faster even with perfect prediction, and every
misprediction makes it slower (\cref{fig:l05-break-even}).  In general the
fully predicated version executes one instruction fewer than the MOVZ/MOVN
version per conditional move, so with $M$ moves its break-even rate is lower
by $M/P$; if the conditional moves test a register that already holds the
condition, they need no compare while full predication still needs one, and
the difference is $(M-1)/P$.  By lowering the break-even rate, full
predication makes if-conversion pay off also for branches that are predicted
too well to justify conditional moves.\margin{ARM's 32-bit A32 instructions
carry a 4-bit condition field, so nearly all of them are predicated; AArch64
(2011) dropped it, keeping conditional select (CSEL).}

\begin{keypoint}[Lesson summary]
  Branch prediction costs nothing when right and tens of instructions when
  wrong; predication executes both directions of a branch and never
  mispredicts, at the price of always doing the unneeded work.  Branches are
  therefore predicted by default, and compilers use predication only for
  small, hard-to-predict if-then-else branches, not for loops, calls, returns
  or large if-then-else branches.  If-conversion removes such a branch by
  turning the control dependence into a data dependence on the condition.
  With ordinary instructions it computes both directions into temporaries and
  selects the results with conditional moves, such as MOVZ and MOVN in MIPS or
  CMOVcc on the condition codes in x86.  If-conversion pays off when the
  mispredictions per execution times the penalty exceed its extra
  instructions (\cref{eq:l05-break-even}); deeper pipelines favour it.  Full
  predication, with a condition on every instruction as in IA-64, avoids the
  temporaries and moves and so lowers the break-even misprediction rate.
\end{keypoint}

\end{document}
